\section{Discussion}
\label{sec:discussion}

\subsection{What the chain implies for evaluating economic agents}

\paragraph{Agreement rate is a poor quality measure in a multi-party market.} A bilateral exchange that completes one player's group raises the risk for every other player, including the recipient. In these games recipients mostly refused such trades or asked for more, and they did so even when the offer paid them. Low acceptance in this population therefore reflects, at least in part, recipients pricing a rival's gain. It does not show that the agents lack a bias toward agreement, because low acceptance can coexist with generous accepted terms, and the cohort has no valuation against which to call a price generous. An evaluation of negotiating agents should report offers by whose position they change and should read counteroffer content, not only accept and reject counts.

\paragraph{Execution is part of economic competence, and the interface helps decide it.} The defaulted Opus builds were affordable and requested, yet the game recorded no development. For a deployed agent the same gap appears whenever a planned allocation must pass through an action language with ordering constraints and a limited retry budget. The relevant measurement separates three things: whether the requested allocation was reachable, whether its encoding was valid, and whether it was executed. Construction counts that merge them measure the interface as much as the agent.

\paragraph{Liquidity is what the rules let you reach.} Displayed options understated finance in all three payable eliminations, and nominal equity overstated it in the G7 lock. An agent or analyst that reasons from either will misjudge solvency. Checking a debt against legally reachable sales and mortgages is a modest computation in this game, and it changed the description of four eliminations.

\paragraph{Terminal transfers belong to the market outcome.} With intact creditor inheritance, a rival's insolvency is a strategic event. It decides who receives the estate, and players argued about this explicitly (G13). Denial purchases held only as long as their holders stayed solvent. Evaluations that score each agent's survival separately miss where its assets went, and in this cohort they went to the eventual winner in 54 of 60 cases.

\paragraph{Conduct audits can start at pressure points.} The selected incidents with the largest stakes occur at solvency and inheritance decisions, and the Arena study reports more misaligned messages under low inventory~\cite{li2026arena}. Both observations come from selected or coded material, so they do not establish a rate. They do suggest where a conduct review of an economic agent should look first: decisions where a party is near insolvency, where a deed would complete a rival's set, or where an estate is about to change hands.

\subsection{The research questions}
\label{sec:questions}

The project began with eight questions about multi-agent behavior and three broader ones. Table~\ref{tab:questions} states what this cohort supports for each and what remains open. Several questions need controlled designs, such as varied opponents, randomized messages or feedback, held-out tasks and independent capability measures, that this observational study does not have. We report that as the state of the evidence.

\begin{table}[p]
\centering\small
\begin{tabularx}{\linewidth}{@{}L{2.6cm}YY@{}}
\toprule
Question & What this cohort supports & What remains open\\
\midrule
Q1 Strategic reasoning & Population evidence that recipients refuse or price a rival's completion, in three configurations through acceptance and in all four through counter content; G11 refuses its own completion because it completes the rival; legal financing tests for every elimination. & Whether choices maximized winning probability; regret against any comparator; optimal policy.\\
Q2 Opponent modeling & Frequent public-capacity deductions about rivals' post-trade cash and liquidation (G2, G11); a few checkable forecasts. & Prediction accuracy and calibration; inference about hidden beliefs; the screenshot's auction-cash example.\\
Q3 Negotiation, cooperation, deception & Complete 920-response census; counter content; initiation and settlement roles; authorized bluffs (G2, G7), an omission with a concealment rationale (G13), a withdrawn assurance (G18), recipient corrections (G11, G19, G6). & Effectiveness against a utility or outside option; agreement bias; causal persuasion. Low acceptance does not establish the absence of agreement bias.\\
Q4 Long-horizon planning, memory & Liquidity at elimination (three payable debts, one locked); realized multi-turn targets (G9); plans that never met an opportunity; memory conditions (Opus's missing notes, formatter defects). & Plan follow-through rates; reserve adequacy before optional spending; the effect of memory.\\
Q5 Adaptation & Local corrections after validator feedback; recurrence of a construction error at later submissions; revised auction ceilings. & Durable learning; adaptation to unfamiliar opponents or repeated exploitation.\\
Q6 Economic decisions, risk & Accessible versus nominal finance; spending and paid receipts by unit; the executability of requested construction. & Ex ante value of a purchase (the \$400 question); risk-adjusted return; incremental return on investment.\\
Q7 Emergent and collective behavior & Concentration of estates in eventual winners; an explicit coordination appeal recognized by its target (G20); blocking positions passing through estates. & Coalition frequency or stability; collusion; effects of communication on collective outcomes.\\
Q8 Evaluation validity & Configured-system description with explicit conditions; replay and ledger reconciliation; configuration and seed-block heterogeneity in the main bargaining result. & Transfer to new opponents, games or models; variance across independent runs; opponent-adjusted comparison.\\
\midrule
A Strategy versus strategic-sounding text & Text and action agree in many incidents (G2, G11, G18), and conflict in others (G12, G19, G16, 139 Grok responses); defaulted plans were executable. & Whether coherent plans reflect strategic competence in general.\\
B General versus multi-agent capability & Not evaluated. Four heterogeneous configurations with unequal conditions cannot estimate this relationship. & An independently measured capability axis across several tiers and families.\\
C Skill versus luck and opponents & Local attribution of particular errors and routes (payable debts, defaults, conflicts, inheritance). Sol wins every game in one seed block; in another, seat 2 wins under all four configurations. & Decomposition of winning into skill, chance and opponent weakness.\\
\bottomrule
\end{tabularx}
\caption{Coverage of the eight research questions and three broader questions. The left column states what the records support at their observed level; the right column states what needs a different design. Question wording follows the project's guidance transcription.}
\label{tab:questions}
\end{table}

Two of these answers deserve a sentence more. For question~A, the cohort gives no general verdict on whether these systems are ``genuinely strategic''. It shows that a complete game contains both kinds of record. In G11 and G2 the stated reason, the applied action and the later game state cohere, while in G12 and G19 the applied action contradicts every text channel in the same response. An evaluation that reads only text, or only actions, would misdescribe some of these decisions. For question~C, the seed-block pattern shows why win totals from this design say little about skill. Shared seeds do not preserve equivalent opportunities once actions change how random draws are consumed, and each configuration always faced the same three opponents.
\FloatBarrier
